\documentclass[11pt,a4paper]{article}

% Packages
\usepackage[utf8]{inputenc}
\usepackage[T1]{fontenc}
\usepackage{amsmath,amssymb,amsthm}
\usepackage{algorithm}
\usepackage{algpseudocode}
\usepackage{booktabs}
\usepackage{hyperref}
\usepackage{graphicx}
\usepackage{geometry}
\usepackage{enumitem}
\usepackage{xcolor}
\usepackage{listings}

\geometry{margin=1in}

% Code listing style
\lstset{
    basicstyle=\ttfamily\small,
    keywordstyle=\color{blue},
    commentstyle=\color{gray},
    stringstyle=\color{red},
    breaklines=true,
    frame=single
}

\title{\textbf{IsoTree: Algorithm Implementation Overview}\\[0.5em]
\large Fair-Cut Forest (FCF) and Isolation Forest Variants}
\author{David\\Based on: ``Revisiting randomized choices in isolation forests'' by David Cortes}
\date{January 2026}

\begin{document}

\maketitle

\begin{abstract}
This document provides a technical overview of the \texttt{isotree} library implementation, focusing on the Fair-Cut Forest (FCF) algorithm proposed in arXiv:2110.13402. We cover the solver framework, computational complexity, parallelism strategies, and internal data structures used in the C++ implementation.
\end{abstract}

\tableofcontents
\newpage

%==============================================================================
\section{Introduction}
%==============================================================================

The \texttt{isotree} library is a high-performance C++ implementation of Isolation Forest and its variants, including:
\begin{itemize}
    \item \textbf{iForest}: Original Isolation Forest (Liu et al., 2008)
    \item \textbf{EIF}: Extended Isolation Forest (Hariri et al., 2019)
    \item \textbf{SCiForest}: Split-Criterion Isolation Forest (Liu et al., 2010)
    \item \textbf{FCF}: Fair-Cut Forest (Cortes, 2021)
    \item \textbf{RRCF}: Robust Random Cut Forest (Guha et al., 2016)
\end{itemize}

The library provides Python, R, and C interfaces, with the core algorithms implemented in C++ for maximum performance.

\textbf{Repository}: \url{https://github.com/david-cortes/isotree}

%==============================================================================
\section{Solver Framework}
%==============================================================================

\subsection{Core Algorithm: Fair-Cut Forest (FCF)}

The FCF algorithm extends the isolation forest concept by using a \textbf{pooled standard deviation gain} criterion for split selection, as opposed to uniformly random splits (iForest) or averaged gain (SCiForest).

\subsubsection{Pooled Gain Criterion}

The key difference from SCiForest lies in the gain calculation:

\begin{equation}
\text{gain}_{\text{pooled}} = \frac{\sigma_{\text{all}} - \frac{n_{\text{left}} \cdot \sigma_{\text{left}} + n_{\text{right}} \cdot \sigma_{\text{right}}}{n_{\text{left}} + n_{\text{right}}}}{\sigma_{\text{all}}}
\end{equation}

Compared to SCiForest's averaged gain:
\begin{equation}
\text{gain}_{\text{avg}} = \frac{\sigma_{\text{all}} - \frac{\sigma_{\text{left}} + \sigma_{\text{right}}}{2}}{\sigma_{\text{all}}}
\end{equation}

The pooled gain tends to produce more \textbf{evenly-divided splits} and better handles \textbf{multimodal distributions} with clustered outliers.

\subsubsection{Algorithm: FairCutTree}

\begin{algorithm}[H]
\caption{FairCutTree}
\begin{algorithmic}[1]
\Require Data points $X_{m \times n}$, number of splitting variables $p$, current depth $d$
\If{$m = 1$}
    \State Set as terminal node with depth $d$
\EndIf
\If{all columns have identical values}
    \State Set as terminal node with depth $d + \mathbb{E}[\text{depth}(m)]$
\EndIf
\State Initialize $\mathbf{z} := \mathbf{0}_m$
\For{$i = 1$ to $p$}
    \State Choose variable $v$ uniformly at random from $[1, n]$ with $>1$ unique value
    \State Draw coefficient $c \sim \mathcal{N}(0, 1)$
    \State Update $\mathbf{z} := \mathbf{z} + c \cdot \frac{\mathbf{y} - \bar{y}}{\sigma_y}$ where $\mathbf{y} = X[:, v]$
\EndFor
\State Find split point $s$ minimizing $\frac{m_{\text{left}} \cdot \sigma_{z_l} + m_{\text{right}} \cdot \sigma_{z_r}}{m_{\text{left}} + m_{\text{right}}}$
\State $X_l = \{x_i \in X \mid z_i \leq s\}$, $X_r = \{x_i \in X \mid z_i > s\}$
\State $\text{Node}_L = \text{FairCutTree}(X_l, p, d+1)$
\State $\text{Node}_R = \text{FairCutTree}(X_r, p, d+1)$
\end{algorithmic}
\end{algorithm}

\subsubsection{Algorithm: Fair-Cut Forest}

\begin{algorithm}[H]
\caption{Fair-Cut Forest}
\begin{algorithmic}[1]
\Require Data $X_{m \times n}$, splitting variables $p$, trees $t$, sample size $s$
\For{$i = 1$ to $t$}
    \State Sample $S$ of $s$ points uniformly at random without replacement
    \State $\text{Tree}_i = \text{FairCutTree}(X_S, p, 0)$
\EndFor
\State Calculate $q = \mathbb{E}[\text{depth}(m)]$
\State \Return Forest $T$ and expected depth $q$
\end{algorithmic}
\end{algorithm}

\subsubsection{Scoring Algorithm}

\begin{algorithm}[H]
\caption{ScorePoint}
\begin{algorithmic}[1]
\Require Point $\mathbf{x}_n$, forest $T$ of $t$ trees, expected depth $q$
\State Initialize $d = 0$
\For{$i = 1$ to $t$}
    \State $d := d + \text{TreeScore}(\mathbf{x}, \text{Root of } T_i)$
\EndFor
\State \Return $2^{-d/t \cdot q^{-1}}$ \Comment{Anomaly score in $[0, 1]$}
\end{algorithmic}
\end{algorithm}

\subsection{Recommended Hyperparameters for FCF}

Based on the paper's experiments:
\begin{itemize}
    \item \textbf{Trees}: 200 (more than iForest's 100 due to higher variability)
    \item \textbf{Sample size}: 32--256 points per tree (smaller for multimodal data)
    \item \textbf{ndim}: 2 variables per split (hyperplane splits)
    \item \textbf{ntry}: 1 (single trial per node)
    \item \textbf{max\_depth}: None (grow until isolation)
    \item \textbf{Coefficients}: $\sim \mathcal{N}(0,1)$
\end{itemize}

%==============================================================================
\section{Time Complexity Analysis}
%==============================================================================

\subsection{Training Complexity}

For a dataset with $m$ observations and $n$ features:

\subsubsection{Single Tree Construction}

\begin{itemize}
    \item \textbf{Sample selection}: $O(s)$ where $s$ is sample size (typically 256)
    \item \textbf{Per-node split evaluation}:
    \begin{itemize}
        \item Linear combination: $O(p \cdot s)$ where $p$ is \texttt{ndim}
        \item Finding optimal split: $O(s \log s)$ for sorting + $O(s)$ for gain calculation
        \item Total per node: $O(s \log s)$
    \end{itemize}
    \item \textbf{Tree depth}: Expected $O(\log s)$ for balanced trees
    \item \textbf{Total nodes}: $O(s)$ in the worst case
\end{itemize}

\textbf{Single tree complexity}: $O(s^2 \log s)$ worst case, $O(s \log^2 s)$ expected

\subsubsection{Forest Construction}

For $t$ trees:
\begin{equation}
T_{\text{train}} = O(t \cdot s \cdot \log^2 s)
\end{equation}

With typical values ($t=200$, $s=256$):
\begin{equation}
T_{\text{train}} \approx O(200 \cdot 256 \cdot 64) = O(3.3 \times 10^6)
\end{equation}

\textbf{Key insight}: Training is \textbf{independent of $m$} when using subsampling, making it highly scalable.

\subsection{Prediction Complexity}

For predicting a single observation:

\begin{itemize}
    \item \textbf{Per tree}: $O(d \cdot p)$ where $d$ is tree depth, $p$ is \texttt{ndim}
    \item \textbf{Expected depth}: $O(\log s)$
    \item \textbf{Total for $t$ trees}: $O(t \cdot p \cdot \log s)$
\end{itemize}

For $k$ observations:
\begin{equation}
T_{\text{predict}} = O(k \cdot t \cdot p \cdot \log s)
\end{equation}

\subsection{Space Complexity}

\subsubsection{Model Storage}

Each node stores:
\begin{itemize}
    \item Split information: $O(p)$ for coefficients, means, std devs
    \item Child pointers: $O(1)$ (indices)
    \item Metadata: $O(1)$
\end{itemize}

Total model size:
\begin{equation}
S_{\text{model}} = O(t \cdot s \cdot p)
\end{equation}

\subsubsection{Comparison Table}

\begin{table}[H]
\centering
\begin{tabular}{@{}lcc@{}}
\toprule
\textbf{Operation} & \textbf{Time Complexity} & \textbf{Space Complexity} \\
\midrule
Training (with subsampling) & $O(t \cdot s \cdot \log^2 s)$ & $O(t \cdot s \cdot p)$ \\
Training (full data) & $O(t \cdot m \cdot \log^2 m)$ & $O(t \cdot m \cdot p)$ \\
Prediction (single) & $O(t \cdot p \cdot \log s)$ & $O(1)$ \\
Prediction (batch of $k$) & $O(k \cdot t \cdot p \cdot \log s)$ & $O(k)$ \\
\bottomrule
\end{tabular}
\caption{Complexity summary for isotree algorithms}
\end{table}

%==============================================================================
\section{Parallelism}
%==============================================================================

\subsection{Multi-threading with OpenMP}

The \texttt{isotree} library uses \textbf{OpenMP} for parallel execution. Key parallelization strategies:

\subsubsection{Tree-level Parallelism}

\begin{lstlisting}[language=C++]
#pragma omp parallel for schedule(dynamic) num_threads(nthreads)
for (size_t tree = 0; tree < ntrees; tree++) {
    // Build tree independently
    build_tree(tree, sample, model);
}
\end{lstlisting}

Each tree is built independently, allowing embarrassingly parallel construction.

\subsubsection{Prediction Parallelism}

\begin{lstlisting}[language=C++]
#pragma omp parallel for schedule(static) num_threads(nthreads)
for (size_t i = 0; i < n_observations; i++) {
    scores[i] = score_observation(X[i], forest);
}
\end{lstlisting}

Predictions for different observations are computed in parallel.

\subsection{Thread Configuration}

From the library documentation:
\begin{itemize}
    \item \textbf{nthreads = -1}: Use all available CPU cores
    \item \textbf{Memory scaling}: Each thread allocates its own buffers
    \item \textbf{Diminishing returns}: Memory bandwidth-bound operations may not scale linearly
\end{itemize}

\textbf{Warning from documentation}:
\begin{quote}
``Adding more threads will not result in a linear speed-up. For some types of data (e.g., large sparse matrices with small sample sizes), adding more threads might result in only a very modest speed-up (e.g., 1.5x faster with 4x more threads).''
\end{quote}

\subsection{SIMD and Architecture Optimizations}

The library supports:
\begin{itemize}
    \item \textbf{-march=native}: Compile-time CPU optimization (AVX, AVX2, etc.)
    \item \textbf{Cache-friendly data layouts}: Flat vector storage instead of pointer-based trees
    \item \textbf{Contiguous memory access}: Improves cache locality during tree traversal
\end{itemize}

%==============================================================================
\section{Data Structures}
%==============================================================================

\subsection{Core Structures (C++)}

The library uses efficient, cache-friendly data structures defined in \texttt{isotree.hpp}:

\subsubsection{IsoTree: Standard Node Structure}

\begin{lstlisting}[language=C++]
typedef struct IsoTree {
    ColType  col_type;           // Numeric, Categorical, or NotUsed
    size_t   col_num;            // Column index for split
    double   num_split;          // Split threshold (numeric)
    std::vector<char> cat_split; // Bitmap for categorical splits
    int      chosen_cat;         // Selected category
    size_t   tree_left;          // Index to left child
    size_t   tree_right;         // Index to right child
    double   pct_tree_left;      // Fraction going left
    double   score;              // Depth/score at node
    double   range_low;          // Range bounds for penalization
    double   range_high;
    double   remainder;          // For distance calculations
} IsoTree;
\end{lstlisting}

\subsubsection{IsoHPlane: Extended/Hyperplane Node}

For extended isolation forest with hyperplane splits:

\begin{lstlisting}[language=C++]
typedef struct IsoHPlane {
    std::vector<size_t>   col_num;   // Multiple columns
    std::vector<ColType>  col_type;  // Type per column
    std::vector<double>   coef;      // Linear combination coeffs
    std::vector<double>   mean;      // Standardization means
    std::vector<double>   fill_val;  // Missing value handling
    std::vector<std::vector<double>> cat_coef; // Categorical coeffs
    
    double   split_point;            // Threshold on combination
    size_t   hplane_left;            // Left child index
    size_t   hplane_right;           // Right child index
    double   score;
    double   range_low, range_high;
    double   remainder;
} IsoHPlane;
\end{lstlisting}

\subsubsection{IsoForest: Complete Model}

\begin{lstlisting}[language=C++]
typedef struct IsoForest {
    std::vector<std::vector<IsoTree>> trees;  // Forest storage
    NewCategAction    new_cat_action;   // Handle new categories
    CategSplit        cat_split_type;   // Categorical strategy
    MissingAction     missing_action;   // Missing value handling
    ScoringMetric     scoring_metric;   // depth, density, etc.
    double            exp_avg_depth;    // Expected isolation depth
    double            exp_avg_sep;      // Expected separation
    size_t            orig_sample_size; // Training sample size
    bool              has_range_penalty;
} IsoForest;
\end{lstlisting}

\subsubsection{ExtIsoForest: Extended Model}

\begin{lstlisting}[language=C++]
typedef struct ExtIsoForest {
    std::vector<std::vector<IsoHPlane>> hplanes;  // Hyperplane trees
    // Same metadata as IsoForest...
} ExtIsoForest;
\end{lstlisting}

\subsection{Design Rationale}

\subsubsection{Flat Vector Storage}

Trees are stored as \textbf{flat vectors} rather than traditional pointer-based nodes:

\begin{itemize}
    \item \textbf{Cache efficiency}: Contiguous memory layout
    \item \textbf{Serialization}: Easy to save/load models
    \item \textbf{Thread safety}: No pointer aliasing issues
    \item \textbf{Memory locality}: Children accessed by index, not pointer
\end{itemize}

\begin{lstlisting}[language=C++]
// Tree traversal by index
size_t current = 0;  // Root
while (!is_terminal(tree[current])) {
    if (x[tree[current].col_num] <= tree[current].num_split)
        current = tree[current].tree_left;
    else
        current = tree[current].tree_right;
}
\end{lstlisting}

\subsubsection{Categorical Handling}

Categorical variables use \textbf{bitmap representation}:

\begin{lstlisting}[language=C++]
std::vector<char> cat_split;  // Bitmap: cat_split[i] = 1 if category i goes left
\end{lstlisting}

This allows $O(1)$ lookup for categorical split decisions.

\subsection{Enumeration Types}

\begin{lstlisting}[language=C++]
typedef enum ColType {
    Numeric = 31, 
    Categorical = 32, 
    NotUsed = 0
} ColType;

typedef enum MissingAction {
    Divide = 21,    // Send to both branches (weighted)
    Impute = 22,    // Fill with learned values  
    Fail = 0        // Raise error
} MissingAction;

typedef enum CategSplit {
    SubSet = 0,        // Subset of categories per branch
    SingleCateg = 41   // One category vs rest
} CategSplit;
\end{lstlisting}

\subsection{Memory Layout Summary}

\begin{table}[H]
\centering
\begin{tabular}{@{}lll@{}}
\toprule
\textbf{Structure} & \textbf{Storage} & \textbf{Purpose} \\
\midrule
\texttt{IsoForest} & \texttt{vector<vector<IsoTree>>} & Standard IF \\
\texttt{ExtIsoForest} & \texttt{vector<vector<IsoHPlane>>} & Extended IF, FCF \\
\texttt{Imputer} & \texttt{vector<vector<ImputeNode>>} & Missing value imputation \\
\texttt{TreesIndexer} & \texttt{vector<SingleTreeIndex>} & Fast distance queries \\
\bottomrule
\end{tabular}
\caption{Main data structure organization}
\end{table}

%==============================================================================
\section{Benchmarks}
%==============================================================================

From the isotree repository benchmarks (100 trees, CovType dataset):

\begin{table}[H]
\centering
\begin{tabular}{@{}lccc@{}}
\toprule
\textbf{Library} & \textbf{Model} & \textbf{256 samples} & \textbf{2048 samples} \\
\midrule
isotree & original & 0.00161s & 0.0848s \\
isotree & extended & 0.00326s & 0.168s \\
scikit-learn & original & 8.30s & 6.89s \\
eif (Python) & extended & 0.16s & 5.06s \\
H2O & original & 9.33s & 14.23s \\
\bottomrule
\end{tabular}
\caption{Training time comparison (lower is better)}
\end{table}

\textbf{Key observation}: isotree is typically 1--2 orders of magnitude faster than other implementations.

%==============================================================================
\section{Conclusion}
%==============================================================================

The \texttt{isotree} library provides a highly optimized implementation of isolation forest variants:

\begin{enumerate}
    \item \textbf{Solver}: Pooled gain criterion (FCF) for better multimodal detection
    \item \textbf{Complexity}: $O(t \cdot s \cdot \log^2 s)$ training, independent of dataset size with subsampling
    \item \textbf{Parallelism}: OpenMP-based tree and prediction parallelism
    \item \textbf{Data structures}: Cache-efficient flat vector storage with index-based traversal
\end{enumerate}

%==============================================================================
\section*{References}
%==============================================================================

\begin{enumerate}
    \item Liu, F.T., Ting, K.M., Zhou, Z.H. (2008). ``Isolation forest.'' ICDM 2008.
    \item Hariri, S., Kind, M.C., Brunner, R.J. (2019). ``Extended Isolation Forest.'' IEEE TKDE.
    \item Liu, F.T., Ting, K.M., Zhou, Z.H. (2010). ``On detecting clustered anomalies using SCiForest.'' ECML-PKDD.
    \item Cortes, D. (2021). ``Revisiting randomized choices in isolation forests.'' arXiv:2110.13402.
    \item Cortes, D. (2021). ``Isolation forests: looking beyond tree depth.'' arXiv:2111.11639.
    \item Guha, S. et al. (2016). ``Robust random cut forest based anomaly detection on streams.'' ICML.
\end{enumerate}

\end{document}
